\documentclass[10pt,a4paper]{amsart}
\usepackage[margin=19mm]{geometry}
\usepackage{amsmath,amssymb,mathtools}
\usepackage[scr=rsfs,cal=euler]{mathalfa}
\usepackage{xfrac}
\usepackage[unicode,hidelinks,bookmarks=false]{hyperref}
\newcommand{\ie}{i.e.\ }
\newcommand{\cf}{cf.\ }
\newcommand{\resp}{respectively\ }
\newcommand{\Subsection}{\subsection}
\providecommand{\nobreakdash}{\nobreak\hbox{-}}
\setlength{\emergencystretch}{2em}
\multlinegap0pt
\usepackage{bm}
\usepackage{bbm}
\usepackage{dsfont}
\usepackage{xspace}
\usepackage{cancel}
\usepackage{mathtools}
\usepackage{amsthm}
\makeatletter
\@ifundefined{theorem}{\newtheorem{theorem}{Theorem}}{}
\@ifundefined{lemma}{\newtheorem{lemma}[theorem]{Lemma}}{}
\makeatother
\newcommand{\V}{\mathcal{V}} 
\newcommand{\cS}{\mathcal{S}} 
\newcommand{\mc}[1]{\mathcal{#1}}
\newcommand{\rgeq}{\trianglerighteq}
\title[A complete characterisation of conditional entropies]{A complete characterisation of conditional entropies}
\author{Roberto Rubboli, Erkka Haapasalo, Marco Tomamichel}
\date{Source: arXiv:2601.23213v1 (CC BY 4.0)}
\begin{document}
\maketitle

\noindent\emph{Source and licence.} A complete characterisation of conditional entropies. Version arXiv:2601.23213v1 (CC BY 4.0), distributed under the Creative Commons Attribution 4.0 International licence, as recorded in the arXiv licence metadata for this exact version and shown on the abstract page https://arxiv.org/abs/2601.23213v1. Full rights notice: licenses/arxiv-2601.23213-CC-BY-4.0.txt.\par\smallskip

\noindent\textit{Editorial scope.} Counted unit: the Lemma whose source label is \texttt{lem: degenerate and zig-zag} in the article (source0.tex lines 911--922, source label \texttt{lem: degenerate and zig-zag}), delivered together with the article's own complete original proof of it (source0.tex lines 923--992, one \texttt{proof} environment of 70 lines). No mathematics is written, repaired or re-derived: statement and proof are the article's own text line for line. Required context retained uncounted: none. Every definition, display and label of those lines is retained unchanged. Omitted from this package and not redistributed: 1--910, 993--3060. Nothing inside a retained excerpt is omitted or altered: every retained body is delivered exactly as the article's lines are written, so no omission receipt of any kind accompanies this package. Citations: the retained excerpts carry no citation command at all, so no bibliography entry of the article is transcribed and no reference is redistributed. Reference adaptation: none was needed and none was made; every reference inside the delivered unit resolves inside it, so no unresolved reference or citation is produced. Printed slips kept literal: no printed word, symbol, hypothesis or number of the counted statement or of its proof was altered. Layout and class: the article's own class is \texttt{revtex4-2} with packages including mathrsfs, graphicx, comment, dsfont, amsmath, amsfonts, amssymb, amsthm, float, hyperref, caption, lipsum, subcaption, mwe; the wrapper is the lane's pinned \texttt{amsart} editorial wrapper at 10pt on A4, which restates the theorem-like environments the retained text needs and adds the article's own macro definitions that the retained text needs (\texttt{\textbackslash V}, \texttt{\textbackslash cS}, \texttt{\textbackslash mc}, \texttt{\textbackslash rgeq}), with extra wrapper packages (bm, bbm, dsfont, xspace, cancel, mathtools). The delivered numbering is the unit's own. Assets: no retained span contains a figure environment, a graphics-inclusion command, a PDF-page inclusion, a table or a TikZ picture; no image file is copied into this package, the package declares no asset dependency and no image file is redistributed with it. Licence: version arXiv:2601.23213v1 of this article is distributed under the Creative Commons Attribution 4.0 International licence, as recorded in the arXiv licence metadata for this exact version and shown on the abstract page https://arxiv.org/abs/2601.23213v1. A full rights notice is delivered at licenses/arxiv-2601.23213-CC-BY-4.0.txt. Characters: every retained excerpt is pure ASCII, so the delivered engine, XeLaTeX with its default Latin Modern text font, reports no missing character.
\begin{lemma}
\label{lem: degenerate and zig-zag}
In the preordered semiring of conditional majorization $S$, the map
\begin{equation}
\|[P_{XY}]\| = \sum_{x\in\mc X}\sum_{y\in\mc Y} P_{XY}(x,y) \,,
\end{equation}
where $P_{XY}$ can be chosen to be any representative of the equivalence class $[P_{XY}]$,
defines a surjective homomorphism $\|\cdot\| : \cS \mapsto \mathbb{R}_+$ with trivial kernel. Moreover, it satisfies
\begin{equation}
[P] \rgeq [Q] \quad \Rightarrow \quad \|[P]\| = \|[Q]\| \quad \Rightarrow\quad  [P] \sim [Q]\,.
\end{equation}
\end{lemma}

\begin{proof}
Clearly, the total weight is a homomorphism since $\|[P]\cdot [Q]\|=\|[P\otimes Q]\|=\|[P]\|\|[Q]\|$ and $\|[P]+[Q]\|=\|[P\oplus Q]\|=\|[P]\|+\|[Q]\|$ for all $[P],[Q]\in\cS$. 

The first direction, namely that $[P] \rgeq [Q] \Rightarrow \|[P]\| = \|[Q]\|$, follows immediately, since conditionally mixing operations preserve the total weight $\|P\|$ of any element $P$. This also shows that the total weight is a degenerate homomorphism since it is invariant under the preorder $\rgeq$.

Let us not turn to the implication $\|[P]\|=\|[Q]\|\Rightarrow [P]\sim [Q]$.  We show that if we define $\|P\|=\|Q\| = c$, then $P\preceq c\succeq Q$ 
and hence $[P]\sim [Q]$. In particular, we show that from the element $[c]$ we can generate both $P$ and $Q$. This is straightforward by considering the matrix with a single non-zero entry $c$ in the top-left corner. Indeed, starting from this matrix, we can construct both $P$ and $Q$. For example, let us focus on $P_{XY} \in \V_{d\times n}$.
Then the steps are the following. Let us denote with $P_{Y}$ the row vectors with components $P_Y(y_i) = \sum_jP_{XY}(x_j,y_i)$. 
First, we create for each outcome $Y$ a column whose first entry is $P_Y(y_i)$ and all other entries are zero. This is done by multiplying on the r.h.s by a row stochastic matrix.
\begin{align}
\begin{pmatrix}
     P_{Y} \\
    0_{(d-1) \times n}   
\end{pmatrix}=
\begin{pmatrix}
    \|P_{XY}\| & 0_{1\times (n-1)} \\
    0_{(d-1)\times 1} & 0_{(d-1)\times (n-1)}
\end{pmatrix} 
\begin{pmatrix}
    \|P_{XY}\|^{-1} P_{Y} \\
    \vdots \\
    \|P_{XY}\|^{-1} P_{Y}    
\end{pmatrix}\,.
\end{align}
Next, we measure the register $Y$ and, conditioned on the outcome $y_i$, we transform the deterministic column (with a single non-zero entry) into a more distributed column $P_{X,Y=y_i}$. Let us denote with $P_{X,Y=y_1}$ the column with components $P_{X,Y=y_1}(x_j)=P_{X,Y}(x_j,y_1)$ and with $P_{X|Y=y_1}$ the column with components $P_{X|Y=y_1}(x_j)=P_{X,Y}(x_j,y_1)/P_Y(y_1)$,
For the first outcome $y_1$, we obtain
\begin{align}
\begin{pmatrix}
     P_{X,Y=y_1} & 0_{d \times (n-1)}  
\end{pmatrix}=
\begin{pmatrix}
    P_{X|Y=y_1} & \star_{1\times (d-1)} 
\end{pmatrix} 
\begin{pmatrix}
     P_{Y} \\
    0_{(d-1) \times n}   
\end{pmatrix}
\begin{pmatrix}
    1 & 0_{1\times (n-1)}\\
    0_{(n-1) \times 1} & 0_{1\times (n-1)} 
\end{pmatrix}
\end{align}
Here, $\star_{1\times (d-1)}$ is any entry that completes the matrix to a doubly stochastic one. By summing over all the outcomes $y_i$ for all $i=1,..,n$, we obtain the desired probability $P_{XY}$. This shows that overall, we have
\begin{equation}
    P_{XY} = \sum_{i=1}^n S^{(i)} P DD^{(i)} \,, 
\end{equation}
where 
\begin{align}
P=\begin{pmatrix}
    \|P_{XY}\| & 0_{1\times (n-1)} \\
    0_{(d-1)\times 1} & 0_{(d-1)\times (n-1)}
\end{pmatrix} \,, \quad 
    D=\begin{pmatrix}
    \|P_{XY}\|^{-1} P_{Y} \\
    \vdots \\
    \|P_{XY}\|^{-1} P_{Y}    
\end{pmatrix} \,,
\end{align}
and for example, for $i=1$, we have
\begin{align}
D^{(1)}=\begin{pmatrix}
    1 & 0_{1\times (n-1)}\\
    0_{(n-1) \times 1} & 0_{1\times (n-1)} 
\end{pmatrix} \,,\quad 
S^{(1)} = \begin{pmatrix}
    P_{X|Y=y_1} & \star_{1\times (d-1)} 
\end{pmatrix}  \,.
\end{align}

\end{proof}

\end{document}
